\begin{textsample}{POS Dim 1 – pb – Score 38.00 – pb/pb\_ef\_27.txt}  \label{ex:f1_pos_128}
5.4 Língua Inglesa 5.4.1 Introdução Olá, colega professor!
Construímos nosso currículo de Língua Inglesa de forma colaborativa, a partir da contribuição de vários professores e estudiosos da área.
Nesse processo, muitas reflexões têm sido importantes.
Em \textbf{um} primeiro momento, teremos espaço para discuti -las. \textbf{Que} tal começarmos pensando \textbf{um} pouco acerca de como a Língua Inglesa era \textbf{vista} e trabalhada em sala de aula e como isso acontece agora?
O processo de \textbf{democratização} da Língua Inglesa foi acentuado a partir do uso da internet e da dinâmica promovida a partir da globalização.
Mydans ( 2007, p. 01 ) \textbf{explicita} \textbf{que} “ Ter \textbf{uma} língua global tem auxiliado a globalização e a globalização tem consolidado a língua global.
Esse processo começou com o domínio dos dois impérios de fala inglesa (... ) e continua até \textbf{hoje} com o novo império virtual da internet ”.
Partindo desses pressupostos, \textbf{salientamos} a importância de se pautar os processos de ensino e aprendizagem de Língua Inglesa em \textbf{uma} perspectiva dinâmica e inovadora.
As necessidades exponenciadas a partir da dinâmica do \textbf{século} XXI emergenciam a recusa a processos de ensino e aprendizagem puramente gramaticais e estáticos.
É importante \textbf{que} os professores \textbf{preocupem} -se em efetivar sua prática \textbf{didático-pedagógica} compreendendo a Língua Inglesa e seu uso \textbf{enquanto} práticas sociais, culturais e, até mesmo, políticas.
Essa realidade predispõe \textbf{que} o ensino e aprendizagem da língua não podem ser \textbf{baseados} nem pautados em concepções estruturalistas, apenas, de forma descontextualizada, mas sim, no uso da língua.
Tal fato não significa \textbf{dizer} \textbf{que} os \textbf{tópicos} gramaticais da língua não são importantes, mas \textbf{que} não devem ser trabalhados de forma isolada, \textbf{baseada} na repetição e sem relação com o uso real e social.
De acordo com Graddol ( 2006 ), após o Old English, o Middle English e o Modern English, passamos a experienciar o quarto período da Língua Inglesa : o inglês como língua global ( ILG ).
Tal fato gera algumas implicações políticas, econômicas, culturais e linguísticas \textbf{que} vão desde o âmbito mundial e refletem também nos aspectos educacionais e na forma em \textbf{que} a Língua Inglesa é abordada nas salas de aulas brasileiras.
PARA REFLETIR...
Kumaravadivelu ( 2012, p. 28 ) contribui \textbf{que} “ \textbf{Uma} língua é melhor aprendida quando o foco não está na língua, quando a atenção dos \textbf{aprendizes} está \textbf{focada} em compreender, \textbf{dizer} e fazer algo com a língua, mesmo se eles não estão explicitamente \textbf{preocupados} com as propriedades formais ”.
A BNCC ( Base Nacional Comum Curricular ), ao abordar a Língua Inglesa, expõe \textbf{que} “ o tratamento dado ao componente prioriza o foco da função social e política do inglês e, nesse sentido, passa a tratá -la em seu status de língua franca ” ( 2017, p. 239 ).
Buscamos aprender inglês através do uso e para o uso. \textbf{Salientamos} \textbf{que}, até então, o ensino das línguas, dinâmicas e vivas, tem se mostrado como \textbf{um} dos pontos fracos da educação básica brasileira.
Portanto, faz -se necessária a superação da dicotomia entre \textbf{teoria} e prática na busca pela melhora e desenvolvimento da qualidade de ensino.
Torna -se, assim, indispensável a efetivação de políticas de formação inicial e continuada, bem como o estabelecimento de metas a serem atingidas em \textbf{termos} de conhecimento e proficiência da língua inglesa, tanto para professores quanto para estudantes.
Os estudantes devem estar no centro do processo de ensino e aprendizagem, desde a elaboração dos documentos normativos até o \textbf{que} lhes foi proposto no “ chão da escola ”.
Esta análise deve ser feita levando -se em consideração seus saberes e suas necessidades, mantendo -se \textbf{um} equilíbrio entre o \textbf{que} eles gostariam e o \textbf{que} eles precisam aprender, propiciando assim \textbf{uma} aprendizagem significativa e enriquecedora.
Assim, torna -se essencial \textbf{que} sejam apresentadas oportunidades para desenvolver habilidades a partir de \textbf{uma} perspectiva de educação linguística consciente, crítica, reflexiva e significativa.
Com vistas à efetivação dessas \textbf{abordagens}, é necessário \textbf{que} os professores lancem mãos de metodologias ativas \textbf{que} envolvam os alunos de maneira potencial e \textbf{que} os coloquem em situação de pró-atividade diante do processo de ensino e aprendizagem em \textbf{que} estão inseridos.
Ensinar e aprender Inglês deixa de ser ouvir e repetir, mas passa a ser, aprender, refletir e produzir.
É necessário \textbf{que} os professores percebam \textbf{que}, para \textbf{que} o aluno aprenda de forma significativa, é essencial \textbf{que} suas práticas estejam pautadas não apenas “ no ensino ” mas também na contínua reflexão acerca do \textbf{que} ensinar, compreendendo \textbf{que} é \textbf{urgente} a redimensão quanto à \textbf{abordagem} dos conteúdos a serem ensinados.
E, acima de tudo, na visão quanto a dimensão entre o ensino e a aprendizagem, compreendendo \textbf{que}, muito além da importância \textbf{que} tem o \textbf{que} se ensina, está a importância naquilo \textbf{que} o aluno aprende, de fato.
Não podemos fugir de tudo \textbf{que} envolve o aprendizado de \textbf{uma} língua, inclusive a gramática, pois todos esses \textbf{tópicos} são essenciais para \textbf{uma} aprendizagem global.
Porém, precisamos dar às \textbf{categorias} gramaticais \textbf{uma} nova dimensão.
O aluno precisa ser levado a pensar e produzir para além de atividades de repetir, decorar frases ou preencher listas/lacunas, sendo capaz de contextualizar e de utilizar a língua no seu uso real, situado.
Os conteúdos deixam de vir em primeiro plano e agora se \textbf{configuram} \textbf{enquanto} meios para \textbf{que} o aluno adquira e aperfeiçoe as habilidades, colocando o uso em evidência e não apenas o “ conteúdo pelo conteúdo ”.
Além disso, o texto da BNCC ( BRASIL, 2017 ) demonstra a necessidade de se pautar o ensino e aprendizagem de Língua Inglesa em \textbf{uma} perspectiva de diversidade de conhecimento de si e do outro, tal como da(s) cultura(s) de si e dos outros.
FALANDO SOBRE DIVERSIDADE...
Schimtz ( 2016 ) \textbf{pontua} acerca do perigo de se efetivar \textbf{um} processo de ensino e aprendizagem \textbf{que} promova atitudes discriminatórias em relação às diversidades do inglês.
Deste modo, é importante considerar a descolonização do conhecimento, aprendendo e refletindo sobre as culturas alvo ( nativas da língua ), mas valorizando e identificando -se com a sua própria cultura, tal como com as demais, fugindo ao hábito da imitação e da suposta superioridade de \textbf{uma} cultura em detrimento das outras, \textbf{que} Anjos ( 2013 ) conceitua como atitudes de supervalorização de \textbf{uma} língua ou cultura.
É \textbf{preciso}, assim, \textbf{que} os professores trabalhem à luz da perspectiva dos usos diversos da língua e não a partir da crença \textbf{que} há \textbf{um} inglês perfeito, com pronúncia \textbf{extremamente} correta e \textbf{que}, assim, torna incorreto e desaprovado o \textbf{que} Barcelos ( 2016 ) chama de “ Brazilian English ”, gerando, inclusive, \textbf{uma} sensação de impotência, frustação e inferioridade.
Nesse contexto, os processos de ensino e aprendizagem \textbf{que} tratam a Língua Inglesa \textbf{enquanto} Língua Global ou Língua Franca, indo além da perspectiva da Língua Estrangeira, horizontalizam as possibilidades de aprendizagem e experiências e se afastam das \textbf{abordagens} \textbf{hegemônicas}.
É necessário \textbf{que} as atividades a serem desenvolvidas nas aulas de Língua Inglesa efetivem \textbf{uma} \textbf{abordagem} de cunho transformador, tornando os alunos aptos a usarem a língua em cenários diversos, de uso real, indo além de estruturas puramente gramaticais.
A formação integral do ser humano \textbf{implica} em \textbf{que} haja entendimento sobre os acontecimentos \textbf{que} nos circundam e \textbf{exige} o nosso \textbf{posicionamento} crítico e atitudinal diante deles.
Desta forma, os professores devem \textbf{fomentar} a curiosidade dos estudantes para \textbf{que} assim eles possam, diante de \textbf{uma} imposição normativa e regulatória, entender o \textbf{que} e por quais razões estão estudando o inglês e quais são as funções sociais e políticas atreladas ao aprendizado deste componente.
Por ser, nos dias \textbf{contemporâneos}, indispensável à formação do cidadão em \textbf{um} mundo onde as \textbf{fronteiras} estão a \textbf{um} clique de computador ou a \textbf{um} toque nos smartphones, os princípios fundamentais do ensino e aprendizagem deste componente devem envolver também o uso das ferramentas digitais, \textbf{instigando} questionamentos e reflexões sobre quais são ou serão suas responsabilidades éticas, sociais e políticas ao lidar com esses conhecimentos, propiciando -lhe \textbf{um} aprendizado crítico e significativo.
Partindo do pressuposto de \textbf{que} lidamos com alunos de \textbf{uma} geração em \textbf{que} as ferramentas digitais são utilizadas não apenas \textbf{enquanto} \textbf{meras} ferramentas, mas redimensionam a maneira em \textbf{que} \textbf{vivem}, interagem, resolvem problemas e, até mesmo, aprendem, consideramos \textbf{que} não podemos insistir no mesmo jeito de ensinar.
Currículos, escolas e processos de ensino e aprendizagem precisam ser reinventados e adaptados a essa nova demanda.
Nesse contexto, a BNCC se desenha a partir de considerações \textbf{que} estão alinhadas com os princípios \textbf{aqui} expostos.
O componente curricular Língua Inglesa para o Ensino Fundamental – Anos Finais, está organizado por eixos, unidades temáticas, objetos de conhecimento e habilidades.
Os eixos estruturantes são cinco : oralidade, leitura, escrita, conhecimentos linguísticos e conhecimentos interculturais.
A \textbf{sistematização} do currículo estadual paraibano foi organizada envolvendo os objetivos de aprendizagem, conteúdos e habilidades, com vistas a facilitar a compreensão e colaborar com o bom desenvolvimento do trabalho.
Além disso, abordaremos também aspectos \textbf{metodológicos} e de avaliação.
VAMOS PENSAR \textbf{UM} POUCO...
Para Rajagopalan ( 2011, p. 65 ), “ A Língua Inglesa \textbf{hoje} é \textbf{uma} língua proteiforme¹.
O \textbf{que} “ rola ” no mundo afora \textbf{hoje} é algo \textbf{que} costumo chamar de “ World English ”, onde falares e sotaques diferentes ( \textbf{que} muitos chamam de “ world englishes ”, no plural ), convivem e, por vezes, se digladiam entre si.
Essa língua não tem pátria, nem está delimitada a \textbf{uma} região geográfica.
É esse fenômeno linguístico \textbf{que} devemos nos esforçar para ensinar e aprender, porque é dele \textbf{que} os \textbf{aprendizes} de \textbf{hoje} vão precisar no futuro bem próximo.

% matched lemmas: abordagem, aprendiz, aqui, basear, categoria, configurar, contemporâneo, democratização, didático-pedagógico, dizer, enquanto, exigir, explicitar, extremamente, focar, fomentar, fronteira, hegemônico, hoje, implicar, instigar, mero, metodológico, pontuar, posicionamento, preciso, preocupar, que, salientar, sistematização, século, teoria, termos, tópico, um, urgente, ver, viver
\end{textsample}
